% Publication source for Written on the Wall II, Conjecture 61.
% Build with: pdflatex -interaction=nonstopmode -halt-on-error PROOF-061.tex
% The companion Markdown note is PROOF-061.md.

% Options for packages loaded elsewhere
\PassOptionsToPackage{unicode}{hyperref}
\PassOptionsToPackage{hyphens}{url}
\PassOptionsToPackage{dvipsnames,svgnames,x11names}{xcolor}
\documentclass[
  11pt,
  letterpaper,
]{article}
\usepackage{xcolor}
\usepackage[margin=1in]{geometry}
\usepackage{amsmath,amssymb}
\setcounter{secnumdepth}{-\maxdimen} % remove section numbering
\usepackage{iftex}
\ifPDFTeX
  \usepackage[T1]{fontenc}
  \usepackage[utf8]{inputenc}
  \usepackage{textcomp} % provide euro and other symbols
\else % if luatex or xetex
  \usepackage{unicode-math} % this also loads fontspec
  \defaultfontfeatures{Scale=MatchLowercase}
  \defaultfontfeatures[\rmfamily]{Ligatures=TeX,Scale=1}
\fi
\usepackage{lmodern}
\ifPDFTeX\else
  % xetex/luatex font selection
\fi
% Use upquote if available, for straight quotes in verbatim environments
\IfFileExists{upquote.sty}{\usepackage{upquote}}{}
\IfFileExists{microtype.sty}{% use microtype if available
  \usepackage[]{microtype}
  \UseMicrotypeSet[protrusion]{basicmath} % disable protrusion for tt fonts
}{}
\makeatletter
\@ifundefined{KOMAClassName}{% if non-KOMA class
  \IfFileExists{parskip.sty}{%
    \usepackage{parskip}
  }{% else
    \setlength{\parindent}{0pt}
    \setlength{\parskip}{6pt plus 2pt minus 1pt}}
}{% if KOMA class
  \KOMAoptions{parskip=half}}
\makeatother
\usepackage{color}
\usepackage{fancyvrb}
\newcommand{\VerbBar}{|}
\newcommand{\VERB}{\Verb[commandchars=\\\{\}]}
\DefineVerbatimEnvironment{Highlighting}{Verbatim}{commandchars=\\\{\}}
% Add ',fontsize=\small' for more characters per line
\newenvironment{Shaded}{}{}
\newcommand{\AlertTok}[1]{\textcolor[rgb]{1.00,0.00,0.00}{\textbf{#1}}}
\newcommand{\AnnotationTok}[1]{\textcolor[rgb]{0.38,0.63,0.69}{\textbf{\textit{#1}}}}
\newcommand{\AttributeTok}[1]{\textcolor[rgb]{0.49,0.56,0.16}{#1}}
\newcommand{\BaseNTok}[1]{\textcolor[rgb]{0.25,0.63,0.44}{#1}}
\newcommand{\BuiltInTok}[1]{\textcolor[rgb]{0.00,0.50,0.00}{#1}}
\newcommand{\CharTok}[1]{\textcolor[rgb]{0.25,0.44,0.63}{#1}}
\newcommand{\CommentTok}[1]{\textcolor[rgb]{0.38,0.63,0.69}{\textit{#1}}}
\newcommand{\CommentVarTok}[1]{\textcolor[rgb]{0.38,0.63,0.69}{\textbf{\textit{#1}}}}
\newcommand{\ConstantTok}[1]{\textcolor[rgb]{0.53,0.00,0.00}{#1}}
\newcommand{\ControlFlowTok}[1]{\textcolor[rgb]{0.00,0.44,0.13}{\textbf{#1}}}
\newcommand{\DataTypeTok}[1]{\textcolor[rgb]{0.56,0.13,0.00}{#1}}
\newcommand{\DecValTok}[1]{\textcolor[rgb]{0.25,0.63,0.44}{#1}}
\newcommand{\DocumentationTok}[1]{\textcolor[rgb]{0.73,0.13,0.13}{\textit{#1}}}
\newcommand{\ErrorTok}[1]{\textcolor[rgb]{1.00,0.00,0.00}{\textbf{#1}}}
\newcommand{\ExtensionTok}[1]{#1}
\newcommand{\FloatTok}[1]{\textcolor[rgb]{0.25,0.63,0.44}{#1}}
\newcommand{\FunctionTok}[1]{\textcolor[rgb]{0.02,0.16,0.49}{#1}}
\newcommand{\ImportTok}[1]{\textcolor[rgb]{0.00,0.50,0.00}{\textbf{#1}}}
\newcommand{\InformationTok}[1]{\textcolor[rgb]{0.38,0.63,0.69}{\textbf{\textit{#1}}}}
\newcommand{\KeywordTok}[1]{\textcolor[rgb]{0.00,0.44,0.13}{\textbf{#1}}}
\newcommand{\NormalTok}[1]{#1}
\newcommand{\OperatorTok}[1]{\textcolor[rgb]{0.40,0.40,0.40}{#1}}
\newcommand{\OtherTok}[1]{\textcolor[rgb]{0.00,0.44,0.13}{#1}}
\newcommand{\PreprocessorTok}[1]{\textcolor[rgb]{0.74,0.48,0.00}{#1}}
\newcommand{\RegionMarkerTok}[1]{#1}
\newcommand{\SpecialCharTok}[1]{\textcolor[rgb]{0.25,0.44,0.63}{#1}}
\newcommand{\SpecialStringTok}[1]{\textcolor[rgb]{0.73,0.40,0.53}{#1}}
\newcommand{\StringTok}[1]{\textcolor[rgb]{0.25,0.44,0.63}{#1}}
\newcommand{\VariableTok}[1]{\textcolor[rgb]{0.10,0.09,0.49}{#1}}
\newcommand{\VerbatimStringTok}[1]{\textcolor[rgb]{0.25,0.44,0.63}{#1}}
\newcommand{\WarningTok}[1]{\textcolor[rgb]{0.38,0.63,0.69}{\textbf{\textit{#1}}}}
\usepackage{longtable,booktabs,array}
\newcounter{none} % for unnumbered tables
\usepackage{calc} % for calculating minipage widths
% Correct order of tables after \paragraph or \subparagraph
\usepackage{etoolbox}
\makeatletter
\patchcmd\longtable{\par}{\if@noskipsec\mbox{}\fi\par}{}{}
\makeatother
% Allow footnotes in longtable head/foot
\IfFileExists{footnotehyper.sty}{\usepackage{footnotehyper}}{\usepackage{footnote}}
\makesavenoteenv{longtable}
\setlength{\emergencystretch}{3em} % prevent overfull lines
\providecommand{\tightlist}{%
  \setlength{\itemsep}{0pt}\setlength{\parskip}{0pt}}
\usepackage{bookmark}
\IfFileExists{xurl.sty}{\usepackage{xurl}}{} % add URL line breaks if available
\urlstyle{same}
\hypersetup{
  pdftitle={Two Structural Reductions for Written on the Wall II{,} Conjecture 61},
  colorlinks=true,
  linkcolor={blue},
  filecolor={Maroon},
  citecolor={Blue},
  urlcolor={blue},
  pdfcreator={LaTeX via pandoc}}

\title{Two Structural Reductions for Written on the Wall II, Conjecture
61}
\author{}
\date{27 July 2026}

\begin{document}
\maketitle

\begin{quote}
\textbf{Status.} Conjecture 61 is not proved here. This note proves the
diameter-at-most-three case, a star-forest certificate theorem, a
one-unit deletion theorem, and a two-connected escape theorem. It then
isolates five explicit open statements: any one of (PI), (DL), or (FA)
would finish the conjecture, (MAXP) implies (PI), and (ESC) would finish
the two-connected case of (DL). Exact computation supports every open
statement on 2,592,586 graph records, but finite verification is not a
proof.
\end{quote}

\section{Abstract}\label{abstract}

For a nontrivial connected graph \(G\), Written on the Wall II,
Conjecture 61 asserts

\[
f(G)\ge
\operatorname{residue}(G)
+\left\lceil\frac{\operatorname{diam}(G)}3\right\rceil,
\tag{C61}
\]

where \(f(G)\) is the maximum order of an induced forest and
\(\operatorname{residue}(G)\) is the Havel--Hakimi residue. We give two
reductions of the conjecture.

The first asks for a distance-three packing \(S\) of order
\(\lceil\operatorname{diam}(G)/3\rceil\) whose deletion leaves an
independent set of order \(\operatorname{residue}(G)\). Such a packing
and independent set induce a disjoint union of stars, proving the
desired bound immediately. A stronger formulation asks that the packing
occur among the vertices removed by a suitable maximum-degree deletion
order.

The second reduction is inductive. If \(v\) is a maximum-degree vertex
and \(C_1,\ldots,C_k\) are the components of \(G-v\), then

\[
\sum_{i=1}^k
\left(
  \operatorname{residue}(C_i)
  +\left\lceil\frac{\operatorname{diam}(C_i)}3\right\rceil
\right)
\ge
\operatorname{residue}(G)
+\left\lceil\frac{\operatorname{diam}(G)}3\right\rceil-1.
\]

Thus the natural component induction can lose at most one vertex. We
state the exact deletion lemma that would remove this last unit and
prove that it implies (C61).

Both reductions are then pushed one step further. In a two-connected
graph, deleting a maximum-degree vertex avoided by some diametral pair
loses nothing at all, so a two-connected counterexample to the deletion
lemma must have every maximum-degree vertex inside every diametral pair,
diameter congruent to \(1\) modulo \(3\), and exact equality in both the
residue and diameter terms after the deletion. Complementing this, a
repair statement is proved: induced forests meeting the component quotas
in which no tree holds two neighbors of the deleted vertex recover the
lost unit whenever they exist. All five open assertions were checked
exactly on every diameter-at-least-four record among the connected
unlabeled graphs through order \(9\), equality-boundary extensions
through order \(11\), and a path-of-gadgets catalogue through order
\(20\).

\section{1. The problem in plain
language}\label{the-problem-in-plain-language}

All graphs in this note are finite, simple, and undirected. The
conjecture compares three quantities. Its current Lean transcription is
recorded as research open {[}6{]}. The \textbf{order} of a graph is its
number of vertices.

\begin{itemize}
\tightlist
\item
  \(f(G)\) is the largest number of vertices that can be retained while
  the induced subgraph on those vertices has no cycle. Every edge
  between retained vertices must remain; edges cannot be pruned
  separately.
\item
  \(\operatorname{diam}(G)\) is the greatest distance between two
  vertices.
\item
  \(\operatorname{residue}(G)\) is obtained from the degree sequence by
  the Havel--Hakimi process: repeatedly remove the largest degree,
  subtract one from that many following degrees, and sort again. When
  only zeros remain, their number is the residue.
\end{itemize}

The residue is a classical lower bound on the independence number
\(\alpha(G)\), so it already certifies an induced forest of
\(\operatorname{residue}(G)\) vertices. Conjecture 61 says that a
connected graph also allows roughly one additional forest vertex for
every three steps of diameter.

Put

\[
r(G)=\operatorname{residue}(G),\qquad
d(G)=\operatorname{diam}(G),\qquad
q(G)=\left\lceil\frac{d(G)}3\right\rceil
\]

and

\[
b(G)=r(G)+q(G),\qquad
\Phi(G)=b(G)-f(G).
\]

The conjecture is the assertion \(\Phi(G)\le0\) for every nontrivial
connected graph. A single graph with \(\Phi(G)>0\) would disprove it.
Graphs with \(\Phi(G)=0\) are the equality boundary.

Complete graphs and stars show that equality is common for unrelated
reasons. If \(n\ge2\), then \(K_n\) has \(r(K_n)=1\), \(d(K_n)=1\), and
\(f(K_n)=2\). The star \(K_{1,n-1}\) has residue \(n-1\), diameter
\(2\), and is itself a forest. Both have \(\Phi=0\).

Long paths with dense decorations are the natural pressure point. The
long spine raises the diameter term, while cliques, odd cycles, and
dense blocks attached along the spine try to make every large induced
subgraph cyclic. The computational constructions below were designed to
push both effects at once.

\section{2. A star-forest certificate}\label{a-star-forest-certificate}

A \textbf{2-packing} is a set \(S\subseteq V(G)\) such that every two
distinct vertices of \(S\) are at distance at least \(3\) in \(G\).

\subsection{Proposition 1 (packing--independence
certificate)}\label{proposition-1-packingindependence-certificate}

Suppose \(G\) contains a 2-packing \(S\) and an independent set
\(I\subseteq V(G)\setminus S\). Then \(G[I\cup S]\) is a disjoint union
of stars and isolated vertices. In particular,

\[
f(G)\ge |I|+|S|.
\]

\subsubsection{Proof}\label{proof}

There are no edges inside \(I\), since \(I\) is independent, and no
edges inside \(S\), since distinct vertices of a 2-packing are not
adjacent. Moreover, each vertex of \(I\) has at most one neighbor in
\(S\). Two such neighbors would be joined by a path of length two
through that vertex, contrary to the definition of a 2-packing.

Every nontrivial component of \(G[I\cup S]\) therefore consists of one
vertex of \(S\) and some of its neighbors in \(I\). It is a star. The
remaining vertices are isolated, so the entire induced subgraph is a
forest. \(\square\)

Proposition 1 turns Conjecture 61 into the following concrete structural
question.

\subsection{Packing--independence lemma (PI,
open)}\label{packingindependence-lemma-pi-open}

Every nontrivial connected graph \(G\) has a 2-packing \(S\) such that

\[
|S|=q(G)
\quad\text{and}\quad
\alpha(G-S)\ge r(G).
\tag{PI}
\]

\subsection{Corollary 2}\label{corollary-2}

If (PI) holds for every nontrivial connected graph, then Conjecture 61
holds.

\subsubsection{Proof}\label{proof-1}

Choose an independent set \(I\) of order \(r(G)\) in \(G-S\), then apply
Proposition 1:

\[
f(G)\ge |I|+|S|=r(G)+q(G)=b(G).
\qquad\square
\]

This reduction explains the divisor \(3\). Vertices chosen three steps
apart cannot acquire a common neighbor in the independent part, so the
combined induced graph has no cycle. Merely choosing mutually
nonadjacent vertices would not be enough.

\section{3. A maximum-degree
formulation}\label{a-maximum-degree-formulation}

The classical \texttt{MAX} algorithm repeatedly deletes a vertex of
maximum degree until the remaining graph is independent. When several
vertices have maximum degree, the result can depend on the tie-breaking.

Griggs and Kleitman proved that every possible independent set left by
\texttt{MAX} has order at least \(r(G)\) {[}2{]}. If \texttt{MAX} stops
with more than \(r(G)\) vertices, we may continue deleting isolated
vertices until exactly \(r(G)\) remain; a zero-degree vertex is then
still a maximum-degree vertex.

This gives a proof-shaped strengthening of (PI).

\subsection{MAX-packing lemma (MAXP,
open)}\label{max-packing-lemma-maxp-open}

For every nontrivial connected graph \(G\), the maximum-degree ties can
be broken so that, after exactly \(|V(G)|-r(G)\) deletions, the deleted
set contains a 2-packing of order \(q(G)\), with distances measured in
the original graph \(G\).

\subsection{Proposition 3}\label{proposition-3}

\[
\text{(MAXP)}\implies\text{(PI)}\implies\text{(C61)}.
\]

\subsubsection{Proof}\label{proof-2}

Run the ordering promised by (MAXP), and let \(D\) be the first
\(|V(G)|-r(G)\) deleted vertices. Its complement \(I\) is independent
and has order \(r(G)\). If \(S\subseteq D\) is the promised 2-packing,
then \(I\subseteq V(G)\setminus S\), so \(\alpha(G-S)\ge r(G)\). This is
(PI), and Corollary 2 finishes the implication. \(\square\)

The existential tie-breaking clause cannot simply be dropped. Among the
21,332 connected unlabeled graphs of order \(9\) and diameter at least
\(4\), all admit at least one successful maximum-degree ordering, but
only 10,604 work for every ordering.

There is also a technical trap in any recursive proof of (MAXP):
deletion can increase distances. A set that is a 2-packing in a later
residual graph need not have been a 2-packing in the original graph. An
induction must therefore remember original distance-two conflicts, not
only the current graph.

\section{4. What can already be
proved}\label{what-can-already-be-proved}

\subsection{Proposition 4 (one vertex beyond
independence)}\label{proposition-4-one-vertex-beyond-independence}

Every nontrivial connected graph satisfies

\[
f(G)\ge\alpha(G)+1.
\]

\subsubsection{Proof}\label{proof-3}

Let \(A\) be a maximum independent set. Since \(G\) is connected and has
at least two vertices, \(A\ne V(G)\); choose \(v\notin A\). The graph
induced by \(A\cup\{v\}\) is a star centered at \(v\), together with
isolated vertices. It is an induced forest of order \(\alpha(G)+1\).
\(\square\)

\subsection{Corollary 5 (diameter at most
three)}\label{corollary-5-diameter-at-most-three}

Conjecture 61 holds whenever \(d(G)\le3\).

\subsubsection{Proof}\label{proof-4}

A nontrivial connected graph has \(1\le d(G)\le3\), hence \(q(G)=1\).
The classical residue bound \(r(G)\le\alpha(G)\) {[}1{]} and Proposition
4 give

\[
f(G)\ge\alpha(G)+1\ge r(G)+q(G).
\qquad\square
\]

The only unresolved case therefore has diameter at least \(4\).

\section{5. The one-unit deletion
theorem}\label{the-one-unit-deletion-theorem}

For a graph \(H\) with connected components \(C_1,\ldots,C_k\), define
the component bound

\[
b^\oplus(H)=
\sum_{i=1}^k
\left(
  r(C_i)+q(C_i)
\right).
\]

For an isolated vertex \(K_1\), we use \(r(K_1)=1\), \(d(K_1)=0\), and
\(q(K_1)=0\), so its contribution to \(b^\oplus\) is \(1\).

\subsection{Lemma 6 (diameter loss under one
deletion)}\label{lemma-6-diameter-loss-under-one-deletion}

If \(G\) is connected, \(v\in V(G)\), and \(C_1,\ldots,C_k\) are the
components of \(G-v\), then

\[
\sum_{i=1}^k q(C_i)\ge q(G)-1.
\tag{1}
\]

\subsubsection{Proof}\label{proof-5}

Let

\[
P=x_0x_1\cdots x_d
\]

be a diametral geodesic of \(G\), so \(d=d(G)\). If \(v\notin V(P)\),
the whole path survives in one component of \(G-v\). Distances cannot
decrease when a vertex is deleted, so that component has diameter at
least \(d\), and the left side of (1) is at least \(q(G)\).

Now suppose \(v=x_i\). If neither endpoint was deleted and \(x_0,x_d\)
lie in one component of \(G-v\), their distance there is at least \(d\),
and the same argument applies. Otherwise, the surviving arm or arms of
\(P-v\) lie in their respective components. Let their edge-lengths be
\(a,b\ge0\), taking the missing arm to have length \(0\) when an
endpoint of \(P\) is deleted. Each nonempty arm remains a geodesic in
its component, and

\[
d\le a+b+2.
\]

Consequently,

\[
\left\lceil\frac a3\right\rceil+
\left\lceil\frac b3\right\rceil
\ge
\left\lceil\frac{a+b+2}{3}\right\rceil-1
\ge
\left\lceil\frac d3\right\rceil-1.
\]

For the first inequality, use \(a\le3\lceil a/3\rceil\) and
\(b\le3\lceil b/3\rceil\). Each existing arm lies in a component whose
diameter is at least its length; a missing arm contributes \(0\). This
proves (1). \(\square\)

\subsection{Lemma 7 (residue under a maximum-degree
deletion)}\label{lemma-7-residue-under-a-maximum-degree-deletion}

Let \(v\) be a maximum-degree vertex of \(G\), and let
\(C_1,\ldots,C_k\) be the components of \(G-v\). Then

\[
\sum_{i=1}^k r(C_i)\ge r(G).
\tag{2}
\]

\subsubsection{Proof}\label{proof-6}

Let \(\pi'\) be the Havel--Hakimi reduction of the degree sequence of
\(G\). It removes the maximum degree and subtracts one from the largest
available terms. Deleting the actual vertex \(v\) instead subtracts one
from the terms corresponding to its neighbors. After sorting, this
actual residual degree sequence dominates \(\pi'\), meaning its initial
partial sums are at least as large: subtracting from any terms other
than the largest ones can only make the sequence more concentrated.
Residue is monotone under this dominance order {[}1{]}, so

\[
r(G-v)\ge r(\pi')=r(G).
\]

The residue of a disjoint union is at most the sum of the residues of
its components {[}3{]}. Therefore

\[
r(G)
\le r(G-v)
\le \sum_{i=1}^k r(C_i).
\qquad\square
\]

\subsection{Theorem 8 (one-unit deletion
theorem)}\label{theorem-8-one-unit-deletion-theorem}

If \(v\) is a maximum-degree vertex of a connected graph \(G\), then

\[
b^\oplus(G-v)\ge b(G)-1.
\tag{3}
\]

\subsubsection{Proof}\label{proof-7}

Add inequalities (1) and (2). \(\square\)

The theorem says that the natural induction is never badly wrong. A
maximum-degree deletion preserves the full residue contribution and can
lose at most one unit from the rounded diameter contribution.

\section{6. The exact inductive finishing
lemma}\label{the-exact-inductive-finishing-lemma}

The one missing unit can be isolated as a second open statement.

\subsection{Deletion lemma (DL, open)}\label{deletion-lemma-dl-open}

Every connected graph \(G\) with \(d(G)\ge4\) has a vertex \(v\) such
that

\[
b^\oplus(G-v)\ge b(G).
\tag{DL}
\]

The vertex in (DL) is not required to have maximum degree, and \(G-v\)
may be disconnected.

\subsection{Theorem 9}\label{theorem-9}

If (DL) holds, then Conjecture 61 holds.

\subsubsection{Proof}\label{proof-8}

Proceed by strong induction on \(|V(G)|\). The case \(d(G)\le3\) is
Corollary 5. Otherwise choose \(v\) as in (DL), and let
\(C_1,\ldots,C_k\) be the components of \(G-v\).

Each nontrivial component has fewer vertices than \(G\), so by induction
it contains an induced forest of order at least \(b(C_i)\). An isolated
component contributes its single vertex. The union of these forests is
an induced forest in \(G-v\), and hence in \(G\), of order at least

\[
\sum_{i=1}^k b(C_i)
=b^\oplus(G-v)
\ge b(G).
\]

This is (C61). \(\square\)

Theorem 8 sharply constrains a smallest counterexample to Conjecture 61.
If \(G\) were such a graph, component induction would give
\(b^\oplus(G-v)\le b(G)-1\) for every vertex \(v\). When \(v\) has
maximum degree, integrality and (3) would therefore force

\[
b^\oplus(G-v)=b(G)-1
\]

for every such \(v\). Both ingredients would have to be tight:

\[
\sum_i r(C_i)=r(G),
\qquad
\sum_i q(C_i)=q(G)-1.
\tag{4}
\]

This reduces the next proof attempt to the equality structure in (4):
show that either another vertex satisfies (DL) or an induced forest
through the deleted vertex recovers the missing unit. The next two
sections carry both alternatives as far as they currently go.

\section{7. The two-connected case}\label{the-two-connected-case}

In the equality structure (4), the diameter term loses its unit only
under a genuine separation. Making this precise closes the two-connected
case up to a single structural question.

\subsection{Lemma 10 (metric loss requires
separation)}\label{lemma-10-metric-loss-requires-separation}

Let \(v\) be a vertex of a connected graph \(G\), and suppose some
diametral pair \(\{x,y\}\) of \(G\) satisfies \(v\notin\{x,y\}\), with
\(x\) and \(y\) in the same component \(C\) of \(G-v\). Then

\[
\sum_{i=1}^k q(C_i)\ge q(G).
\]

\subsubsection{Proof}\label{proof-9}

Deleting a vertex cannot decrease distances, so the distance between
\(x\) and \(y\) inside \(C\) is at least \(d(G)\). Hence
\(d(C)\ge d(G)\) and \(q(C)\ge q(G)\). \(\square\)

This is the first case in the proof of Lemma 6, restated because its
contrapositive is what matters here: losing the diameter unit forces the
deleted vertex to separate, or to belong to, every diametral pair. No
vertex separates a two-connected graph.

\subsection{Theorem 11 (two-connected
escape)}\label{theorem-11-two-connected-escape}

Let \(G\) be two-connected with \(d(G)\ge4\), and let \(v\) be a
maximum-degree vertex avoided by some diametral pair. Then \(G-v\) is
connected and

\[
b(G-v)\ge b(G).
\]

In particular, \(v\) satisfies (DL).

\subsubsection{Proof}\label{proof-10}

Since \(G\) is two-connected, \(G-v\) is connected, so the avoided
diametral pair lies in the single component \(G-v\) and Lemma 10 gives
\(q(G-v)\ge q(G)\). Lemma 7 gives \(r(G-v)\ge r(G)\). Adding the two
inequalities proves the claim. \(\square\)

\subsection{Corollary 12 (structure of a trapped
counterexample)}\label{corollary-12-structure-of-a-trapped-counterexample}

Suppose \(G\) is two-connected with \(d(G)\ge4\) and no vertex of \(G\)
satisfies (DL). Then:

\begin{enumerate}
\def\labelenumi{\arabic{enumi}.}
\tightlist
\item
  every maximum-degree vertex lies in every diametral pair; consequently
  \(G\) has at most two maximum-degree vertices, and if it has exactly
  two, the diametral pair is unique;
\item
  \(d(G)\equiv1\pmod3\);
\item
  for every maximum-degree vertex \(v\), both \(d(G-v)=d(G)-1\) and
  \(r(G-v)=r(G)\) hold exactly;
\item
  for every diametral pair \(\{v,y\}\) with \(v\) of maximum degree,
  \(v\) is the only vertex of \(G\) at distance \(d(G)\) from \(y\).
\end{enumerate}

\subsubsection{Proof}\label{proof-11}

Claim 1 is Theorem 11 read contrapositively; a diametral pair has two
elements, and two distinct maximum-degree vertices can both lie in every
pair only if every pair is exactly those two vertices.

For claim 4, if some \(z\ne v\) had \(\operatorname{dist}(z,y)=d(G)\),
then \(\{z,y\}\) would be a diametral pair avoiding \(v\), against claim
1.

For claims 2 and 3, fix a maximum-degree \(v\) and a diametral pair
\(\{v,y\}\). The graph \(G-v\) is connected and retains the arm of a
diametral geodesic from \(y\), so \(d(G-v)\ge d(G)-1\), while Lemma 7
gives \(r(G-v)\ge r(G)\). If \(d(G)\not\equiv1\pmod3\), then
\(\lceil(d(G)-1)/3\rceil=q(G)\) and \(b(G-v)\ge b(G)\), contradicting
the absence of a (DL) vertex. So \(d(G)\equiv1\pmod3\), and
\(b(G-v)\le b(G)-1\) forces equality throughout: \(r(G-v)=r(G)\) and
\(q(G-v)=q(G)-1\). Since \(d(G)\equiv1\pmod3\) gives
\(3(q(G)-1)=d(G)-1\), the ceiling equality caps \(d(G-v)\le d(G)-1\), so
\(d(G-v)=d(G)-1\) exactly. \(\square\)

The remaining question is whether trapped graphs exist at all.

\subsection{Escape lemma (ESC, open)}\label{escape-lemma-esc-open}

Every two-connected graph with \(d(G)\ge4\) has a maximum-degree vertex
avoided by some diametral pair.

By Theorem 11, (ESC) implies (DL) for every two-connected graph. All
53,096 two-connected diameter-at-least-four records in the four corpora
satisfy (ESC); no trapped graph has been seen (Section 9.4). Two
cautions attach to that evidence. First, coverage: in a two-connected
graph every intermediate breadth-first level from any vertex has at
least two vertices, since a singleton level would be a cut vertex
separating what lies beyond it; hence \(|V(G)|\ge2\,d(G)\), and the
exhaustive order-at-most-\(9\) slice tests (ESC) only at \(d(G)=4\).
Evidence at larger diameters comes from the structured families alone.
Second, a trapped graph would not by itself refute (DL) or Conjecture
61: Corollary 12 would still leave its rigid residual structure open to
a direct argument.

\section{8. Forests through the deleted
vertex}\label{forests-through-the-deleted-vertex}

The second alternative in the equality analysis of (4) keeps the deleted
vertex in the forest.

\subsection{Proposition 13 (free-attachment
repair)}\label{proposition-13-free-attachment-repair}

Let \(v\) be a maximum-degree vertex of a connected graph \(G\) with
\(d(G)\ge4\), and let \(C_1,\ldots,C_k\) be the components of \(G-v\).
Suppose each \(C_i\) contains an induced forest \(F_i\) with
\(|F_i|\ge b(C_i)\) such that no tree of \(C_i[F_i]\) contains two
neighbors of \(v\). Then

\[
f(G)\ge b(G).
\]

\subsubsection{Proof}\label{proof-12}

The union \(F=F_1\cup\cdots\cup F_k\) is an induced forest of \(G-v\).
Any cycle in \(G[F\cup\{v\}]\) would have to pass through \(v\), and its
remainder would be a path in the forest \(G[F]\) joining two neighbors
of \(v\). A path in a forest stays inside one tree, so some tree would
contain two neighbors of \(v\), contrary to hypothesis. Adding \(v\)
merges the trees it touches into one and creates no cycle, so
\(G[F\cup\{v\}]\) is an induced forest of order at least

\[
1+\sum_{i=1}^k b(C_i)
=1+b^\oplus(G-v)
\ge b(G)
\]

by Theorem 8. \(\square\)

\subsection{Free-attachment lemma (FA,
open)}\label{free-attachment-lemma-fa-open}

Every connected graph \(G\) with \(d(G)\ge4\) has a maximum-degree
vertex \(v\) admitting forests as in Proposition 13.

\subsection{Proposition 14}\label{proposition-14}

If (FA) holds, then Conjecture 61 holds.

\subsubsection{Proof}\label{proof-13}

Corollary 5 covers \(d(G)\le3\); Proposition 13 covers \(d(G)\ge4\)
directly. \(\square\)

The hypotheses of (FA) cannot simply be dropped. If a component is a
single edge with both endpoints adjacent to \(v\), the only forest
meeting the quota \(b(K_2)=2\) is the edge itself, and its single tree
holds both neighbors. That configuration is a triangle and has diameter
\(1\); the open content of (FA) is that diameter at least \(4\) and
maximum degree leave enough room. The star-forest witness of Section 2
is the natural candidate for the \(F_i\): stars are hard to enter twice,
which is one more sign that the packing reduction and the deletion
reduction aim at the same object.

At every maximum-degree deletion in the corpora that is one-unit tight
--- \(b^\oplus(G-v)=b(G)-1\), the only case Theorem 8 leaves open, since
otherwise \(v\) itself satisfies (DL) --- the constrained quota forests
exist: 58,809 configurations, no exceptions (Section 9.4).

\section{9. Exact computational
evidence}\label{exact-computational-evidence}

The computation had two roles: search directly for \(\Phi(G)>0\), and
test the proposed structural lemmas without replacing them by
heuristics.

\subsection{9.1 Evaluator validation}\label{evaluator-validation}

The optimized C evaluator was compared with separate exhaustive oracles
on all 33,867 labeled graphs through order \(6\):

\begin{itemize}
\tightlist
\item
  manual Havel--Hakimi reduction against the production residue routine;
\item
  Floyd--Warshall against repeated-BFS diameter;
\item
  all vertex subsets against the cycle-branching induced-forest
  optimizer;
\item
  graph6 encoding and decoding round trips.
\end{itemize}

Complete graphs, paths, and stars were also checked through order
\(20\). An independent Python program re-evaluated reported boundary
graphs by descending subset enumeration.

\subsection{9.2 Counterexample search}\label{counterexample-search}

Four complementary searches were used.

\begin{enumerate}
\def\labelenumi{\arabic{enumi}.}
\tightlist
\item
  \textbf{Exhaustive small graphs.} Every connected unlabeled graph
  through order \(9\) was generated. The order-\(9\) count of 261,080
  agrees with McKay's published table {[}4{]}.
\item
  \textbf{Path-of-gadgets constructions.} The catalogue joins cliques,
  odd cycles, complete bipartite blocks, barbells, articulation blocks,
  mixed blocks, and layered paths along thin necks, through order
  \(20\).
\item
  \textbf{Boundary extensions.} Every nonempty neighborhood for a new
  vertex was attached to each of the 1,231 order-\(9\) equality graphs.
  Leaves and true or false twins were then attached to the 49,521
  surviving order-\(10\) equality records.
\item
  \textbf{Direct optimization.} Edge additions and deletions preserving
  connectivity hill-climbed \(\Phi\) and the smoother score
  \(3(r(G)-f(G))+d(G)\). About 7.5 million edge toggles were considered
  at orders \(10,12,13,16,\) and \(18\), with every connected candidate
  evaluated exactly.
\end{enumerate}

No graph with \(\Phi(G)>0\) was found. This is a bounded null result,
not a proof of (C61).

\subsection{9.3 Structural-lemma checks}\label{structural-lemma-checks}

\href{proof61.c}{\texttt{proof61.c}} enumerates 2-packings and solves
the remaining independent-set decision problem exactly.
\href{maxine61.c}{\texttt{maxine61.c}} uses dynamic programming over
remaining-vertex masks, explores every allowed maximum-degree tie, and
inspects the frontier at exactly \(r(G)\) remaining vertices.

{\def\LTcaptype{none} % do not increment counter
\begin{longtable}[]{@{}
  >{\raggedright\arraybackslash}p{(\linewidth - 12\tabcolsep) * \real{0.1111}}
  >{\raggedleft\arraybackslash}p{(\linewidth - 12\tabcolsep) * \real{0.1481}}
  >{\raggedleft\arraybackslash}p{(\linewidth - 12\tabcolsep) * \real{0.1481}}
  >{\raggedleft\arraybackslash}p{(\linewidth - 12\tabcolsep) * \real{0.1481}}
  >{\raggedleft\arraybackslash}p{(\linewidth - 12\tabcolsep) * \real{0.1481}}
  >{\raggedleft\arraybackslash}p{(\linewidth - 12\tabcolsep) * \real{0.1481}}
  >{\raggedleft\arraybackslash}p{(\linewidth - 12\tabcolsep) * \real{0.1481}}@{}}
\toprule\noalign{}
\begin{minipage}[b]{\linewidth}\raggedright
corpus
\end{minipage} & \begin{minipage}[b]{\linewidth}\raggedleft
records
\end{minipage} & \begin{minipage}[b]{\linewidth}\raggedleft
\(d\le3\)
\end{minipage} & \begin{minipage}[b]{\linewidth}\raggedleft
\(d\ge4\)
\end{minipage} & \begin{minipage}[b]{\linewidth}\raggedleft
(PI)
\end{minipage} & \begin{minipage}[b]{\linewidth}\raggedleft
some MAX order
\end{minipage} & \begin{minipage}[b]{\linewidth}\raggedleft
every MAX order
\end{minipage} \\
\midrule\noalign{}
\endhead
\bottomrule\noalign{}
\endlastfoot
connected unlabeled graphs of order 9 & 261,080 & 239,748 & 21,332 &
21,332 & 21,332 & 10,604 \\
all one-vertex extensions of order-9 equality graphs & 629,041 & 509,354
& 119,687 & 119,687 & 119,687 & 37,579 \\
leaf/twin extensions of surviving order-10 equality records & 1,342,198
& 962,218 & 379,980 & 379,980 & 379,980 & 111,752 \\
path-of-gadgets catalogue through order 20 & 360,267 & 376 & 359,891 &
359,891 & 359,891 & 213,501 \\
\textbf{total evaluated records} & \textbf{2,592,586} &
\textbf{1,711,696} & \textbf{880,890} & \textbf{880,890} &
\textbf{880,890} & \textbf{373,436} \\
\end{longtable}
}

The extension and family rows contain isomorphic duplicates. The totals
therefore count exact evaluations, not distinct isomorphism classes. A
separate component-bound scan found a vertex satisfying (DL) in all
880,890 records of diameter at least \(4\).

The deletion reduction gives a more detailed picture at order \(9\). Of
the 21,332 graphs not covered by Corollary 5:

\begin{itemize}
\tightlist
\item
  15,568 have a successful maximum-degree deletion that leaves a
  connected graph;
\item
  21,136 have some successful connected deletion;
\item
  the remaining 196 have no successful connected deletion, but all 196
  have a cut-vertex deletion satisfying (DL).
\end{itemize}

Thus no order-\(9\) graph is induction-critical. Thirty-one of the final
196 graphs also lie on the equality boundary \(\Phi=0\).

\subsection{9.4 Endgame checks}\label{endgame-checks}

\href{finish61.c}{\texttt{finish61.c}} ran three further exact checks
across the same four corpora. Its internal deletion census independently
reproduced the order-\(9\) numbers above.

\textbf{Lemma 7 stress test.} All 5,246,681 maximum-degree deletions
across the 2,592,586 records satisfied \(\sum_i r(C_i)\ge r(G)\), with
zero violations. This exercises the dominance step of Lemma 7 --- the
one step of the proved material resting on a monotonicity property cited
from {[}1{]} --- on every available record.

\textbf{(ESC) and (FA) census.} Every two-connected record of diameter
at least \(4\) has a maximum-degree vertex avoided by some diametral
pair, and the conclusion \(b(G-v)\ge b(G)\) of Theorem 11 was
re-verified by direct computation at each such vertex, with no failures.
Every one-unit-tight maximum-degree deletion admits the constrained
quota forests of Proposition 13; the per-component search was exhaustive
over vertex subsets, with no records skipped.

{\def\LTcaptype{none} % do not increment counter
\begin{longtable}[]{@{}
  >{\raggedright\arraybackslash}p{(\linewidth - 8\tabcolsep) * \real{0.1579}}
  >{\raggedleft\arraybackslash}p{(\linewidth - 8\tabcolsep) * \real{0.2105}}
  >{\raggedleft\arraybackslash}p{(\linewidth - 8\tabcolsep) * \real{0.2105}}
  >{\raggedleft\arraybackslash}p{(\linewidth - 8\tabcolsep) * \real{0.2105}}
  >{\raggedleft\arraybackslash}p{(\linewidth - 8\tabcolsep) * \real{0.2105}}@{}}
\toprule\noalign{}
\begin{minipage}[b]{\linewidth}\raggedright
corpus
\end{minipage} & \begin{minipage}[b]{\linewidth}\raggedleft
2-connected \(d\ge4\)
\end{minipage} & \begin{minipage}[b]{\linewidth}\raggedleft
trapped
\end{minipage} & \begin{minipage}[b]{\linewidth}\raggedleft
tight deletions
\end{minipage} & \begin{minipage}[b]{\linewidth}\raggedleft
repaired
\end{minipage} \\
\midrule\noalign{}
\endhead
\bottomrule\noalign{}
\endlastfoot
connected unlabeled graphs of order 9 & 1,944 & 0 & 854 & 854 \\
all one-vertex extensions of order-9 equality graphs & 18,037 & 0 &
7,113 & 7,113 \\
leaf/twin extensions of surviving order-10 equality records & 30,380 & 0
& 28,936 & 28,936 \\
path-of-gadgets catalogue through order 20 & 2,735 & 0 & 21,906 &
21,906 \\
\textbf{total} & \textbf{53,096} & \textbf{0} & \textbf{58,809} &
\textbf{58,809} \\
\end{longtable}
}

\section{10. False strengthenings}\label{false-strengthenings}

Exact data rules out three tempting shortcuts.

\begin{enumerate}
\def\labelenumi{\arabic{enumi}.}
\item
  Replacing residue by independence is false:

  \[
  f(G)\ge\alpha(G)+q(G)
  \]

  fails on 539 connected graphs of order \(9\).
\item
  Letting \(\rho_2(G)\) be the maximum order of a 2-packing, the
  stronger inequality

  \[
  f(G)\ge r(G)+\rho_2(G)
  \]

  fails on 114 order-\(9\) graphs with diameter at least \(4\). A
  maximum packing need not leave enough room for the independent part.
\item
  Replacing \(\alpha(G-S)\ge r(G)\) in (PI) by \(r(G-S)\ge r(G)\) is too
  strong. Ten connected order-\(9\) graphs have no target-size packing
  satisfying the residue condition, although each has a packing
  satisfying the independence condition in (PI).
\end{enumerate}

These failures locate the useful level of generality. The target-size
packing and the independent set left after its deletion survive; the
more obvious stronger statements do not.

\section{11. What remains}\label{what-remains}

There are now three independent finishing routes,

\[
\text{(MAXP)}\implies\text{(PI)}\implies\text{(C61)},
\qquad
\text{(DL)}\implies\text{(C61)},
\qquad
\text{(FA)}\implies\text{(C61)},
\]

and one partial route: by Theorem 11, (ESC) implies the two-connected
case of (DL).

The two-connected case is the most constrained. Corollary 12 pins any
counterexample there to diameter congruent to \(1\) modulo \(3\), at
most two maximum-degree vertices locked inside every diametral pair, and
exact equality in both deletion terms. Settling (ESC) is the natural
next step. A trapped graph, if one exists, has at least \(2\,d(G)\ge14\)
vertices when \(d(G)\ge7\), so it lies beyond the exhaustive corpus;
targeted constructions --- a high-degree fan feeding a long ladder, at
diameters \(7\) and \(10\) --- would either produce one or sharpen the
evidence. For \(d(G)=4\) the exhaustive data covers only orders \(8\)
and \(9\).

On the block-tree side, (FA) is the statement that still lacks a
mechanism. The 58,809 successful repairs say the constrained forests
always exist in the tested range; they do not say why. The mod-\(3\)
arithmetic of Lemma 6 favors cut vertices at positions congruent to
\(1\) modulo \(3\) along a diametral geodesic --- the same spacing as
the packing skeleton of Section 2 --- but such a vertex must also
preserve the residue sum, and residue is unprotected for deletions below
maximum degree.

A counterexample to (DL), (PI), (MAXP), (ESC), or (FA) would not
necessarily be a counterexample to Conjecture 61; it would only rule out
that proof route. Further unconstrained edge mutation is unlikely to
resolve the question; the equality classification that Section 6 called
for has now been carried into Sections 7 and 8, and what remains of it
is exactly (ESC) and (FA).

The related WOWII Conjecture 17 says that the maximum order of an
induced bipartite subgraph is at least
\(\alpha(G)+\lceil d(G)/3\rceil\); it is listed as resolved on the
historical WOWII page {[}5{]}. Its diameter term has the same spacing
geometry, but bipartiteness permits even cycles. Proposition 1 uses
distance three to obtain the stronger conclusion that the witness is a
forest.

\section{12. Reproduction}\label{reproduction}

The implementation is in this directory. All build products are written
to \texttt{/tmp}.

\begin{Shaded}
\begin{Highlighting}[]
\BuiltInTok{cd}\NormalTok{ problems/wow2}
\FunctionTok{make}\NormalTok{ test}
\FunctionTok{make}\NormalTok{ check{-}counts{-}9}
\FunctionTok{make}\NormalTok{ exhaustive{-}61}
\FunctionTok{make}\NormalTok{ families{-}61}
\FunctionTok{make}\NormalTok{ check{-}critical{-}61}
\FunctionTok{make}\NormalTok{ check{-}maxine{-}61}
\FunctionTok{make}\NormalTok{ check{-}finish{-}61}
\end{Highlighting}
\end{Shaded}

The direct order-\(9\) packing scan is:

\begin{Shaded}
\begin{Highlighting}[]
\ExtensionTok{/tmp/math{-}eval{-}wow2/wow2{-}generate}\NormalTok{ 9 }\DataTypeTok{\textbackslash{}}
  \KeywordTok{|} \ExtensionTok{/tmp/math{-}eval{-}wow2/wow2{-}proof61} \AttributeTok{{-}{-}scan{-}packing}
\end{Highlighting}
\end{Shaded}

The exact evaluator and graph routines are in
\href{graph.c}{\texttt{graph.c}} and
\href{conjectures.c}{\texttt{conjectures.c}}. The endgame checks of
Section 9.4 are implemented in \href{finish61.c}{\texttt{finish61.c}};
any graph6 stream can be piped through the resulting
\texttt{wow2-finish61}. The independent verifier is
\href{verify.py}{\texttt{verify.py}}. The complete bounded-search log,
including graph6 records on the equality boundary, is in
\href{RESULT-061.md}{\texttt{RESULT-061.md}}.

\section{References}\label{references}

\begin{enumerate}
\def\labelenumi{\arabic{enumi}.}
\tightlist
\item
  O. Favaron, M. Mahéo, and J.-F. Saclé, ``On the residue of a graph,''
  \emph{Journal of Graph Theory} \textbf{15} (1991), 39--64.
  \url{https://doi.org/10.1002/jgt.3190150107}
\item
  J. R. Griggs and D. J. Kleitman, ``Independence and the Havel--Hakimi
  residue,'' \emph{Discrete Mathematics} \textbf{127} (1994), 209--212.
  \url{https://doi.org/10.1016/0012-365X(92)00479-B}
\item
  D. Amos, R. Davila, and R. Pepper, ``On the \(k\)-residue of disjoint
  unions of graphs with applications to \(k\)-independence,''
  \emph{Discrete Mathematics} \textbf{321} (2014), 24--34.
  \url{https://doi.org/10.1016/j.disc.2013.12.013}
\item
  B. D. McKay, ``Combinatorial data: graphs.''
  \url{https://users.cecs.anu.edu.au/~bdm/data/graphs.html}
\item
  E. DeLaViña, ``Written on the Wall II: Resolved Conjectures.''
  \url{https://cms.dt.uh.edu/faculty/delavinae/research/wowII/resolvedT.htm}
\item
  Google DeepMind, ``FormalConjectures/WrittenOnTheWallII/
  GraphConjecture61.lean.''
  \url{https://github.com/google-deepmind/formal-conjectures/blob/main/FormalConjectures/WrittenOnTheWallII/GraphConjecture61.lean}
\end{enumerate}

\end{document}
